\begin{figure}[!t]
\centering
\begin{tikzpicture}[
  every node/.style={font=\footnotesize},
  box/.style={draw, rounded corners=2pt, fill=#1,
              minimum width=5.0cm, minimum height=0.70cm,
              text width=4.8cm, align=center},
  wide/.style={draw, rounded corners=2pt, fill=#1,
               minimum width=8.0cm, minimum height=0.70cm,
               text width=7.8cm, align=center},
  arr/.style={-{Stealth[length=5pt]}, thick},
  bnd/.style={draw=gray, dashed, thick, rounded corners=4pt},
  lbl/.style={font=\scriptsize\itshape, text=gray},
  abi/.style={font=\scriptsize\itshape, text=gray, fill=white},
]

%% Zig
\node[wide=green!10] (prog) at (0, 0) {%
  \texttt{allocator.alloc(u8, n)}\\
  \scriptsize Zig program, given a \texttt{std.mem.Allocator}};

\node[wide=green!20] (vtable) at (0,-1.6) {%
  \texttt{std.mem.Allocator} vtable: \texttt{.alloc}, \texttt{.free}\\
  \scriptsize indirect call, unless the allocator is \texttt{comptime}};

\node[box=green!30] (zposix) at (-3.0,-3.4) {%
  \texttt{posix} allocator\\
  \scriptsize calls \texttt{mmap(null, len, \ldots)}};

\node[box=green!30] (zdirect) at (3.0,-3.4) {%
  \texttt{direct} allocator\\
  \scriptsize calls \texttt{includeos\_zig\_alloc(len, alignment)}};

%% C linkage
\node[box=blue!10] (musl) at (-3.0,-5.9) {%
  \texttt{mmap()} $\rightarrow$ \texttt{syscall\_SYS\_mmap()}\\
  \scriptsize musl, patched to call IncludeOS};

\node[box=blue!10] (shim) at (3.0,-5.9) {%
  \texttt{includeos\_zig\_alloc()}\\
  \scriptsize \texttt{try} \ldots\ \texttt{catch}: exceptions become \texttt{nullptr}};

%% C++
\node[box=blue!20] (sysmmap) at (-3.0,-8.4) {%
  \texttt{sys\_mmap()} $\rightarrow$ \texttt{\_\_sys\_mmap()} $\rightarrow$ \texttt{kalloc()}\\
  \scriptsize typed POSIX layer (\cref{fig:mmap-callchain})};

\node[wide=orange!30] (raw) at (0,-10.4) {%
  \texttt{os::mem::raw\_allocator().allocate()}\\
  \scriptsize \texttt{mem\_resource}, a \texttt{buddy\_resource} by default};

%% arrows
\draw[arr] (prog) -- (vtable);
\draw[arr] (vtable.south) -- ++(0,-0.4) -| (zposix.north);
\draw[arr] (vtable.south) -- ++(0,-0.4) -| (zdirect.north);
\draw[arr] (zposix) -- (musl);
\draw[arr] (zdirect) -- (shim);
\draw[arr] (musl) -- (sysmmap);
\draw[arr] (sysmmap.south) -- ++(0,-0.45) -| ([xshift=-1.5cm]raw.north);
\draw[arr] (shim.south) -- ++(0,-2.95) -| ([xshift=1.5cm]raw.north);

%% layers
\draw[bnd] (-5.9, 0.75) rectangle (5.9,-4.30);
\node[lbl, rotate=90, anchor=south] at (-5.9,-1.78) {Zig};

\draw[bnd] (-5.9,-5.05) rectangle (5.9,-6.75);
\node[lbl, rotate=90, anchor=south] at (-5.9,-5.90) {C};

\draw[bnd] (-5.9,-7.50) rectangle (5.9,-11.15);
\node[lbl, rotate=90, anchor=south] at (-5.9,-9.33) {\CC{} (IncludeOS)};

\node[abi] at (0,-4.68) {C ABI (\texttt{extern "C"})};
\node[abi] at (0,-7.13) {\CC{} ABI (name mangling, exceptions)};

\end{tikzpicture}
\caption{The two allocation paths a Zig program can take on IncludeOS. Both
  start from the same \texttt{std.mem.Allocator} interface and leave Zig
  through the C ABI. The \texttt{posix} allocator calls \texttt{mmap} in musl,
  which goes through the whole POSIX layer. The \texttt{direct} allocator
  calls a small shim, which goes straight to \texttt{raw\_allocator()} and
  stops any \CC{} exception from reaching Zig. Both end up in the same
  allocator.}
\label{fig:zig-allocator-path}
\end{figure}
